\begin{exo}{}
	On considère une fonction $g$ dont la représentation graphique est donnée ci-dessous.

	\begin{center}
		\begin{tikzpicture}[x=0.35cm, y=0.5cm, >=Triangle]
			% Grille secondaire / principale
			\draw[gray!50, thin, step=1] (-11,-5) grid (12,7);

			% Axes
			\draw[->, thick] (-11,0) -- (12,0) node[above] {$x$};
			\draw[->, thick] (0,-5) -- (0,7) node[left] {$y$};

			% Graduations principales et étiquettes
			\foreach \x in {-10, -5, 1, 5, 11} {
					\draw (\x,2pt) -- (\x,-2pt);
				}
			\node[below, font=\footnotesize] at (-10,0) {$-10$};
			\node[below, font=\footnotesize] at (-5,0) {$-5$};
			\node[below left, font=\footnotesize] at (0,0) {$0$};
			\node[below, font=\footnotesize] at (1,0) {$1$};
			\node[below, font=\footnotesize] at (5,0) {$5$};
			\node[below, font=\footnotesize] at (11,0) {$11$};

			\foreach \y in {-4, 1, 5} {
					\draw (2pt,\y) -- (-2pt,\y);
				}
			\node[left, font=\footnotesize] at (0,-4) {$-4$};
			\node[left, font=\footnotesize] at (0,1) {$1$};
			\node[left, font=\footnotesize] at (0,5) {$5$};

			% Courbe de la fonction g
			\draw[line width=1.3pt, blue, samples=100] (-10,0)
			.. controls (-8.9,5.23) and (-6.9,4.73) .. (-5,2)
			.. controls (-3,-0.2) and (-1,0) .. (0,0)
			.. controls (0.5,0) and (1,0) .. (1.5,-0.2)
			.. controls (3,-0.8) and (6,-4) .. (8,-4)
			.. controls (9.5,-4) and (10.5,1) .. (11,6);

			% Points remarquables (extrémités)
			\fill[blue] (8,-4) circle (2pt);
			\fill[blue] (0,0) circle (2pt);
			\fill[blue] (-7.5,4) circle (2pt);
			\fill[blue] (-10,0) circle (2pt);
			\fill[blue] (11,6) circle (2pt);

			% Nom de la courbe
			% \node[blue, above right] at (-5.25,1.75) {$\Cg$};
			\node[blue] at (-10.3,2) {$\Cg$};
		\end{tikzpicture}
	\end{center}

	\begin{enumerate}
		\item Décrire les variations de $g$ sur son ensemble de définition $\mathcal{D}$ et préciser son minimum et son maximum.
	\end{enumerate}
\end{exo}
